% Part 3 opening chapter carrying "Landauer-Based Thermodynamic Model for the Minimum Quasiparticle Density in Al/AlOx/Al Josephson Junction Transmon
% Qubits" (14 Apr 2026), read in full 2026-10-02 (854 lines), with the corrected model. Numbers: docs/verification/scripts/verify_xqp.py.
% Check and referee note: docs/verification/particle/XQP_CHECK.md, XQP_REFEREE_NOTE.md.
\chapter{The cost of holding a qubit}\label{ch:xqp}

\section*{The place}
IAM's Law (Chapter~\ref{ch:iams_law}) prices every irreversible record at $k_BT\ln2$ per bit at the nearest encoding surface. A superconducting qubit
is a device built to \emph{avoid} records: it holds a superposition for as long as it can. This chapter asks what the law says that holding costs, and
where the cost should appear in a measurement.

\section{The anomaly}
In the best-isolated Al/AlOx/Al transmons the density of broken Cooper pairs (quasiparticles) does not fall below
\begin{equation}x_{\rm qp}=n_{\rm qp}/n_{\rm cp}\sim10^{-7},\qquad n_{\rm cp}=2\nu_0\Delta\approx4\times10^6\,\mu{\rm m}^{-3},\end{equation}
that is $n_{\rm qp}\approx0.4$ per $\mu$m$^3$. In thermal equilibrium at 15--20\,mK there should be no more than one quasiparticle in a volume the size
of the Earth. Known external sources have been reduced one by one --- stray infrared light, environmental radioactivity, cosmic rays --- and the floor
remains. The same qubit operated at the surface and deep underground at Gran Sasso, with thirty times fewer muons, shows the same $T_1$. The excess
quasiparticles have a thermalised energy distribution, sitting at the gap edge, which a continuous slow source gives and a bursty external one does not.
Their tunnelling statistics are Poissonian. And the tunnelling rate is sensitive to the quality of the oxide barrier, which contains high-transmission
sites with a reduced gap. Together these point to a source inside the device, at the junction.

\section{The proposal}
A qubit holds its state against a bath of two-level fluctuators (TLS) in the oxide and at the metal--oxide interfaces. A TLS flip kicks the condensate
phase by $\delta\phi\sim g/\omega_q\sim10^{-3}$; the condensate restores its phase in $\hbar/\Delta\approx3.6$\,fs, ten orders of magnitude faster than
the next flip ($\tau_{\rm TLS}\sim1$--$100\,\mu$s). After restoration the condensate is the same whichever way the TLS flipped: the information about
the flip has been discarded. That is an erasure, and by IAM's Law it is paid in heat. Each active site therefore produces an irreversible event at rate
$1/\tau_{\rm TLS}$, whether or not the qubit is being used. The separation of time scales makes the events independent, which gives the observed
Poissonian statistics.

\section{From events to quasiparticles}
Three facts of superconductor physics fix how events become a density.
\paragraph{Pair breaking is quantised.} A quasiparticle exists only once a Cooper pair is broken, which costs $2\Delta$ ($364\,\mu$eV in Al) and yields
two quasiparticles. Landauer's principle sets the \emph{minimum} heat of an erasure, $\Delta\ln2=126\,\mu$eV if the bath sits at the gap scale;
heat below $2\Delta$ goes into phonons that cannot break pairs. An erasure event that produces quasiparticles therefore breaks a pair: yield $2$ per
event.
\paragraph{Quasiparticles spread.} In aluminium they diffuse $\sqrt{D\tau_{\rm qp}}\approx250$--$800\,\mu$m within a lifetime $\tau_{\rm qp}\sim100\,\mu$s,
across the whole island. What a qubit measures is the island average.
\paragraph{The density is per pair near the gap.} $n_{\rm cp}=2\nu_0\Delta$ counts the electrons within $\Delta$ of the Fermi surface, not all
conduction electrons ($n_e/2=9\times10^{10}\,\mu$m$^{-3}$, which is $22{,}600$ times larger).
With $N$ active sites on an island of volume $V$, the steady state is
\begin{equation}x_{\rm qp}=\frac{2N\,\tau_{\rm qp}}{\tau_{\rm TLS}\,n_{\rm cp}\,V}.\label{eq:xqp}\end{equation}

\section{What it requires}
\begin{center}\small\begin{tabular}{lcc}\toprule
island volume $V$ & active sites for $x_{\rm qp}=10^{-7}$ & $x_{\rm qp}$ from one site\\\midrule
$10^3\,\mu$m$^3$ & 60 & $1.7\times10^{-9}$\\ $10^4\,\mu$m$^3$ & 600 & $1.7\times10^{-10}$\\ $10^5\,\mu$m$^3$ & 6{,}000 & $1.7\times10^{-11}$\\\bottomrule
\end{tabular}\end{center}
($\tau_{\rm TLS}=30\,\mu$s, $\tau_{\rm qp}=100\,\mu$s.) A few pinholes cannot hold the floor. Hundreds to thousands of active fluctuators can, and
TLS spectroscopy maps them mostly on the junction leads, which is where the source must be. The number $N$ is a property of each device and has to be
measured; nothing else in Eq.~\eqref{eq:xqp} is free.

\section{The test}
On one device, measure the five quantities $x_{\rm qp}$ (parity switching), $\tau_{\rm qp}$, $\tau_{\rm TLS}$ (1/$f$ noise), the active fluctuator count
$N$ (TLS spectroscopy) and the island volume $V$. The law predicts
\begin{equation}\frac{x_{\rm qp}\,n_{\rm cp}\,V\,\tau_{\rm TLS}}{N\,\tau_{\rm qp}}=2.\end{equation}
Across devices, the floor should scale as $N/(\tau_{\rm TLS}V)$: better barriers (longer $\tau_{\rm TLS}$, fewer fluctuators) lower it; better
shielding does not. The proposal fails if the floor does not follow $N/V$, or if it keeps falling with shielding after external sources are gone.

\section{The coherence ceiling}
Quasiparticle relaxation is $\Gamma_1=x_{\rm qp}\,\omega_q$ to within a factor of order one, so a floor of $10^{-7}$ at $5$\,GHz caps $T_1$ near $0.3$\,ms,
the plateau of the best aluminium transmons. If the floor is the cost of holding the state, it is lowered by changing the barrier and its
fluctuators, not by shielding the device.

\section{Status}
Stands: an endogenous floor from erasures at the TLS rate, consistent with the underground result, the thermalised distribution and the Poissonian
statistics. Derived here: the density in the field's units, with quantised pair breaking and diffusion over the island (Eq.~\ref{eq:xqp}). Open: the
active-site count $N$ for a measured device, and whether each erasure event breaks a pair (yield $2$) or only some do, which would lower the yield and
raise the required $N$.
